\section*{Chapter \ref{ch:m2:interest-rate-swaps} --- Interest-Rate Swaps}

\begin{solution}{exo:m2:interest-rate-swaps:1}
$100\,000\,000 \times 0.0405 \times 365/360 = \text{USD}~4\,106\,250$. It receives SOFR
compounded over the same 365 days, actual/360, on USD~100 million, paid two
business days after the period ends.
\end{solution}

\begin{solution}{exo:m2:interest-rate-swaps:2}
Receiving fixed is long duration, like owning a thirty-year bond financed at
the overnight rate: it loses when rates rise, as its liabilities (which it is
hedging) fall in value too.
\end{solution}

\begin{solution}{exo:m2:interest-rate-swaps:3}
$4.05 - 4.20 = -15$ basis points. Before 2008, when \omterm{def:m2:interest-rate-swaps:swap}{floating legs} paid an
interbank rate with bank credit risk, a negative spread would have meant that
the market priced banks' term borrowing below the government's, which did not
happen; \omterm{def:m2:interest-rate-swaps:spread}{swap spreads} were positive.
\end{solution}

\begin{solution}{exo:m2:interest-rate-swaps:4}
$N A \times 10^{-4}$ is the sensitivity to a change in the swap's own fixed rate
holding the curve's discount factors; bumping the par input and rebuilding also
changes the discount factors that value the \omterm{def:m2:interest-rate-swaps:swap}{fixed leg}, so the two differ by a
second-order term (USD~16 on 82\,000). Hedge with the rebuilt, bucketed number:
it is what the hedging instruments' prices actually do.
\end{solution}

\begin{solution}{exo:m2:interest-rate-swaps:5}
A payer 5y5y gains USD~82\,651 per basis point on the ten-year input and loses
USD~45\,394 on the five-year: it is long the ten-year rate and short the five-year,
since it is the ten-year swap minus the five-year swap. Hedge: receive fixed on
a ten-year par swap with the same ten-year DV01 (USD~100.5 million, since a
ten-year par swap of USD~100 million has 82\,262 on that pillar) and pay fixed
on a five-year par swap with 45\,394 (USD~100.3 million, at 45\,278 per 100
million), leaving only the small residuals at the other pillars.
\end{solution}

\begin{solution}{exo:m2:interest-rate-swaps:6}
Gross notional USD~300 million; each dealer pays fixed on one swap and receives
on another of identical terms, so its net risk is zero. Compression tears up
all three swaps: gross notional falls to zero, and so do the margin and capital
the three trades consumed, with no change in anyone's risk.
\end{solution}

\begin{solution}{exo:m2:interest-rate-swaps:7}
Four-year par rate 3.8425\%, eight-year 3.9743\%, ten-year discount factor
0.666773. The intermediate par rates depend on the interpolation; the inputs do
not.
\end{solution}

\begin{solution}{exo:m2:interest-rate-swaps:8}
An overnight-rate swap has almost no credit risk in its \omterm{def:m2:interest-rate-swaps:swap}{floating leg}, and its
counterparty risk is removed by clearing: its rate is not a bank's borrowing
cost. A negative spread says that receiving fixed on a swap is cheaper to hold
than owning a Treasury, which must be financed on a costly balance sheet, and
that long-dated demand for receiving (pension funds) exceeds what dealers will
absorb. It is a price of balance sheet and of demand, not of credit.
\end{solution}

\begin{solution}{pb:m2:interest-rate-swaps:1}
\textbf{1.} 4.05\%; all-in $4.05 + 1.50 = 5.55\%$.
\textbf{2.} $+\text{USD}~164\,474$ per basis point: the payer gains when rates rise.
\textbf{3.} Zero at the mid-market par rate.
\textbf{4.} The ten-year pillar only.
\textbf{5.} Dollar swaps start two business days after the trade: 25 September
2026 was a Friday.
\textbf{6.} $+\text{USD}~8.02$ million.
\textbf{7.} $-\text{USD}~8.44$ million.
\textbf{8.} A \omterm{def:m2:interest-rate-swaps:payer}{payer swap} is short a fixed-rate bond and long a floating one; it
is short convexity, so it gains less on a rise than it loses on an equal fall.
\textbf{9.} Shocks $-20.0$, $-16.1$, $-12.2$, $-4.4$, $+3.3$ and $+15.0$ basis points at
one, two, three, five, seven and ten years; value $+\text{USD}~2.46$ million.
\textbf{10.} Its risk sits entirely on the ten-year pillar, which moved 15 basis
points: $15 \times 164\,474 = 2.47$ million, less a small convexity term.
\textbf{11.} No: the swap's loss is offset by the lower value of the fixed-rate
debt it has synthetically created; its interest cost is still 5.55\%. The loss
is the value of having fixed at 4.05\% when the market would now fix lower.
\textbf{12.} To book the swap's changes in value against the hedged item's, so
that reported earnings do not swing with rates when the economic position is
hedged; treasurers care because volatile earnings are costly to explain.
\textbf{13.} Two basis points for ten years on its DV01: about $2 \times 164\,474 =
\text{USD}~328\,947$ of present value.
\textbf{14.} Each is exposed to the other's default when the swap has value to
it; managed by collateral under a credit support annex, by central clearing
for dealers, and for corporate clients by credit lines and charges in the
price.
\textbf{15.} The interest of years six to ten, which would float again.
\textbf{16.} It keeps the bank loan's flexibility and pricing, can be adjusted or
unwound separately, and fixes the rate without a bond issue's documentation
and investor base.
\textbf{17.} By receiving fixed in the interdealer market or on a SEF, or with
Treasury futures and bonds, in the same DV01 and pillar, then managing the
residual curve risk in its book.
\textbf{18.} The company has floating debt and wants to fix what it pays: it pays
fixed. A pension fund has long fixed liabilities and wants assets that gain
when rates fall: it receives fixed.
\textbf{19.} Named result: \emph{the corporate hedge} fixes at 4.05\% (5.55\% all in),
with a DV01 of USD~164\,474 per basis point, and is worth USD~2.46 million after the
steepening.
\textbf{20.} A fixed rate for a floating one, on a notional that never changes
hands.
\end{solution}

\begin{solution}{iq:m2:interest-rate-swaps:1}
An exchange of fixed-rate interest for floating-rate interest on a notional
amount, over a schedule. A company with floating-rate debt pays fixed to lock
its cost; one with fixed debt may receive fixed to benefit from falling rates;
an investor uses it to add or remove duration without buying or selling bonds.
\iqlookfor{the exchange of cash flows and a concrete use.}
\end{solution}

\begin{solution}{iq:m2:interest-rate-swaps:2}
Compounded overnight interest over a period is replicated by rolling a deposit,
so in a single-curve world each period's payment is worth $P(t_{i-1}) -
P(t_i)$, and the whole leg $P(t_0) - P(t_n)$. With separate discount and
projection curves, project the compounded rate from the forward curve and
discount on the collateral curve.
\iqlookfor{the telescoping argument, and awareness of multi-curve pricing.}
\end{solution}

\begin{solution}{iq:m2:interest-rate-swaps:3}
Choose instruments (deposits or futures at the front, swaps beyond), pillars
at their maturities, and an interpolation (of log discount factors, zero rates
or forwards); solve for each pillar in turn so that its instrument reprices,
with a root finder, or all at once. Test: every input reprices to within a
tolerance; forwards have no spurious oscillation; bumping one input changes
the curve locally; the curve is stable to the order of inputs; published
curves are reproduced.
\iqlookfor{the choices (instruments, interpolation) and a test list.}
\end{solution}

\begin{solution}{iq:m2:interest-rate-swaps:4}
Owning a Treasury ties up balance sheet and must be financed in repo; receiving
fixed on a swap needs only margin. Long-dated duration demand, notably from
pension funds, is met more cheaply in swaps, and dealers' balance sheets are
too costly to arbitrage the gap away. The \omterm{def:m2:interest-rate-swaps:swap}{floating leg}, once LIBOR and now SOFR
or an overnight rate, carries little credit risk, so nothing pushes swap rates
above government yields.
\iqlookfor{balance sheet and demand, not credit.}
\end{solution}

\begin{solution}{iq:m2:interest-rate-swaps:5}
Receiving a 5y5y is long the ten-year rate's decline and short the five-year's:
it is the ten-year receiver minus the five-year receiver. Hedge by paying fixed
on a ten-year par swap and receiving on a five-year par swap in the
corresponding DV01s; the residual is exposure to the shape between the
pillars, which depends on interpolation.
\iqlookfor{the decomposition into two par swaps.}
\end{solution}

\begin{solution}{iq:m2:interest-rate-swaps:6}
Non-linearity: the parallel bump is not the sum of the local bumps for a
convex position (second-order terms); interpolation that couples pillars;
bumping par rates in one run and zero rates in another; different curve
rebuild settings or tolerances between runs; instruments outside the pillar
set (extrapolation); a discount curve bumped in one computation and not in
the other; numerical noise from a loose root-finder tolerance.
\iqlookfor{second-order effects versus implementation inconsistencies.}
\end{solution}
